\documentclass[11pt]{article}
\usepackage[T1]{fontenc}
\usepackage[margin=1in]{geometry}
\usepackage[colorlinks=true,linkcolor=blue,urlcolor=blue]{hyperref}
\usepackage{array}
\usepackage{booktabs}
\usepackage{parskip}

\title{\textsf{absfigure}: User Guide}
\author{Christian Holz}
\date{2026-07-18 v0.3}

\newcommand{\pkg}{\textsf{absfigure}}
\newcommand{\env}[1]{\texttt{#1}}
\newcommand{\key}[1]{\texttt{#1}}

\begin{document}
\maketitle

\section{What the package does}

\pkg{} lets you place figures and tables on an explicit page, at the top or
bottom, and in a specific column when the document is two-column. In other
words, it supports requests such as:

\begin{itemize}
\item put this figure at the top of page 3, left column
\item put this table at the bottom of page 6, right column
\item put this wide figure across both columns on page 2
\end{itemize}

Unlike absolute positioning packages that place content outside the float
system, \pkg{} still uses \LaTeX{} float machinery. Captions, labels,
references, list-of-figures, and list-of-tables continue to work normally, and
regular \env{figure}/\env{table} floats can still appear alongside absolute
ones.

\section{Provided environments}

The package provides four environments:

\begin{center}
\begin{tabular}{ll}
\toprule
Environment & Meaning \\
\midrule
\env{absfigure} & single-column figure \\
\env{absfigure*} & double-column figure \\
\env{abstable} & single-column table \\
\env{abstable*} & double-column table \\
\bottomrule
\end{tabular}
\end{center}

All four environments use the same optional key syntax:

\begin{verbatim}
\begin{absfigure}[page=3,pos=t,col=2]
  ...
\end{absfigure}
\end{verbatim}

\section{Quick start}

\begin{verbatim}
\documentclass[twocolumn]{article}
\usepackage{graphicx}
\usepackage{absfigure}

\begin{document}

\begin{absfigure}[page=3,pos=t,col=1]
  \centering
  \includegraphics[width=\columnwidth]{fig.pdf}
  \caption{Top of page 3, left column}
  \label{fig:left}
\end{absfigure}

\begin{absfigure*}[page=4,pos=b]
  \centering
  \includegraphics[width=\linewidth]{wide.pdf}
  \caption{Bottom of page 4, full width}
  \label{fig:wide}
\end{absfigure*}

\begin{abstable}[pos=b,col=r]
  \centering
  \caption{Bottom of the current page, right column}
  \label{tab:results}
  \begin{tabular}{lc}
    \toprule
    Method & Score \\
    \midrule
    Ours & 94.2 \\
    \bottomrule
  \end{tabular}
\end{abstable}

\end{document}
\end{verbatim}

\section{Keys}

\begin{center}
\begin{tabular}{>{\ttfamily}l>{\ttfamily}l>{\ttfamily}l p{2.7in}}
\toprule
Key & Values & Default & Meaning \\
\midrule
page & integer, current, end & current & target page number \\
pos & t, top, b, bottom & t & top or bottom placement \\
col & 1, l, left, 2, r, right & 1 & target column; ignored for starred environments \\
\bottomrule
\end{tabular}
\end{center}

\subsection*{\key{page}}

\begin{itemize}
\item Omit \key{page}, or write \key{page=current}, to target the page being
  built where the environment appears.
\item Use \key{page=<number>} to target a specific page.
\item Use \key{page=end} to append a final page containing the float.
\end{itemize}

\subsection*{\key{pos}}

\begin{itemize}
\item \key{pos=t} or \key{pos=top} places the float at the top of the target
  page or column.
\item \key{pos=b} or \key{pos=bottom} places the float at the bottom of the
  target page or column.
\end{itemize}

\subsection*{\key{col}}

\begin{itemize}
\item \key{col=1}, \key{col=l}, or \key{col=left} selects the left column.
\item \key{col=2}, \key{col=r}, or \key{col=right} selects the right column.
\item For \env{absfigure*} and \env{abstable*}, the package ignores
  \key{col} because the float spans both columns.
\end{itemize}

\section{Common patterns}

\subsection*{Single-column figure on a specific page}

\begin{verbatim}
\begin{absfigure}[page=5,pos=t,col=2]
  \centering
  \includegraphics[width=\columnwidth]{example.pdf}
  \caption{Top of page 5, right column}
\end{absfigure}
\end{verbatim}

\subsection*{Wide figure across both columns}

\begin{verbatim}
\begin{absfigure*}[page=2,pos=t]
  \centering
  \includegraphics[width=\linewidth]{overview.pdf}
  \caption{Spanning figure on page 2}
\end{absfigure*}
\end{verbatim}

Inside \env{absfigure*} and \env{abstable*}, \verb|\linewidth| is the full text
width, as with ordinary \env{figure*} and \env{table*}.

\subsection*{Table on the current page}

\begin{verbatim}
\begin{abstable}[pos=b,col=1]
  \centering
  \caption{Bottom of the current page}
  \begin{tabular}{lc}
    ...
  \end{tabular}
\end{abstable}
\end{verbatim}

\subsection*{Appendix-style final page}

\begin{verbatim}
\begin{absfigure*}[page=end,pos=t]
  \centering
  \includegraphics[width=\linewidth]{teaser.pdf}
  \caption{Placed on a final appended page}
\end{absfigure*}
\end{verbatim}

\section{How placement behaves}

\subsection*{Ordering}

If several absolute floats target the same slot, meaning the same combination
of \key{page}, \key{col}, and \key{pos}, they are placed in source order.

Pages processed by the \textsf{balance} package are supported. Absolute floats
remain attached to their requested physical columns after balancing.

\subsection*{Interaction with \textsf{balance} and ordinary floats}

The \textsf{balance} package leaves balancing enabled after \verb|\balance|;
it remains active until \verb|\nobalance|. If \verb|\balance| is used before a
forced page boundary to balance only the preceding material, disable it
immediately after the boundary:

\begin{verbatim}
\balance
% final material on the page to balance
\newpage       % or \clearpage
\nobalance
\end{verbatim}

Leaving \verb|\balance| active across later pages can make ordinary
\env{figure} and \env{table} floats appear out of source order. This behavior
comes from \textsf{balance}, which builds and splits combined column material;
it is not caused by the absolute-float environments. \pkg{} deliberately does
not redefine the standard float environments or insert implicit float
barriers.

The usual \LaTeX{} rules also continue to apply to ordinary floats. In
particular, a single-column \verb|[H]| float can overtake deferred
\env{figure*} or \env{table*} floats. When source order across such a semantic
boundary matters, load \textsf{placeins} and insert \verb|\FloatBarrier|
explicitly at that boundary.

\subsection*{Reruns}

If a float targets a page that has already been shipped out in the current run,
\LaTeX{} cannot go back and modify that earlier page immediately. In that case,
\pkg{} saves the float definition and replays it on the next run.

When this happens, the log includes a warning telling you that placement is not
yet final and that you should rerun \LaTeX{}.

This is normal in cases such as:

\begin{itemize}
\item a float defined on page 6 that targets page 4
\item major edits that shift page breaks
\item multiple absolute floats competing for the same constrained slot
\end{itemize}

\subsection*{What happens when a float does not fit}

The package does not crop or overlap floats. If the requested page, position,
and column do not have enough room, the float may remain unplaced and the log
will report that fact. Typical reasons are:

\begin{itemize}
\item the float is taller than the available space
\item too many floats are competing for the same slot
\item the requested placement is structurally impossible
\end{itemize}

\section{What the package loads}

For normal use, add only:

\begin{verbatim}
\usepackage{absfigure}
\end{verbatim}

\pkg{} loads its own direct dependencies. Users do not need to preload the
packages it requires internally.

You still need to load content-specific packages when your float body uses
them, for example:

\begin{itemize}
\item \texttt{graphicx} for \verb|\includegraphics|
\item \texttt{subcaption} for \env{subfigure}
\item \texttt{booktabs}, \texttt{tabularx}, or similar table packages
\end{itemize}

\section{Installation}

\subsection*{Project-local installation}

Copy \texttt{absfigure.sty} into your project directory and load it with
\verb|\usepackage{absfigure}|.

\subsection*{Maintainer / CTAN source layout}

The canonical source files are:

\begin{itemize}
\item \texttt{source/latex/absfigure/absfigure.dtx}
\item \texttt{source/latex/absfigure/absfigure.ins}
\end{itemize}

To regenerate the runtime package from the documented source:

\begin{verbatim}
latex source/latex/absfigure/absfigure.ins
\end{verbatim}

\section{Where to look next}

The repository also includes:

\begin{itemize}
\item \texttt{README.md} for a short overview
\item \texttt{doc/latex/absfigure/absfigure-test.tex} for a demo document
\item \texttt{source/latex/absfigure/absfigure.dtx} for the package source
\end{itemize}

\end{document}
